\chapter*{如何使用本书}
\addcontentsline{toc}{chapter}{如何使用本书}

本书按经典教材方式组织，但它同时也是一本实验书。每章都应按下面顺序学习：

\begin{enumerate}
  \item 阅读本章目标和起始问题。
  \item 读正文，手工推导关键例子。
  \item 运行章节中的命令。
  \item 修改例子，观察输出如何变化。
  \item 查看汇编、IR、二进制或性能数据。
  \item 完成实验。
  \item 完成作业并写报告。
  \item 对照验收标准检查自己是否真正掌握。
\end{enumerate}

\section*{每章的正确读法}

本册每章都可以按“问题、模型、证据、报告”四步读。

\begin{enumerate}
  \item \textbf{问题：}先明确本章回答什么问题。例如“CPU 如何执行指令”不是问 CPU 有多快，而是问指令如何改变寄存器、内存和控制流。
  \item \textbf{模型：}建立足够简单但不错误的模型。例如先把内存看成按字节寻址的数组，再逐步叠加对象、对齐、cache line 和虚拟地址。
  \item \textbf{证据：}每个模型都要找到工具证据。源码对应编译命令，汇编对应 \code{objdump}，运行时状态对应 GDB，性能结论对应原始样本和环境记录。
  \item \textbf{报告：}把观察写下来。没有报告，就很难发现自己把语言层、ABI 层、硬件层和测量层混在了一起。
\end{enumerate}

读正文时不要只划重点。底层知识的重点不是一句定义，而是你能否把定义用于一个真实例子。例如“指针是地址”这句话不够；你要能说明指针值、指针类型、对象生命周期、地址计算、load 指令和虚拟地址之间分别是什么关系。

\section*{三条主线}

本书有三条主线。

\subsection*{第一条：抽象到实现}

从 C/C++ 语言抽象出发，追到编译器、汇编、机器码、CPU 执行和操作系统机制。

\subsection*{第二条：模型到证据}

先建立模型，再用工具观察。比如先理解寄存器、栈和调用约定，再用 \code{gdb}、\code{objdump}、\code{readelf} 和 benchmark 数据验证。

\subsection*{第三条：基础到工程}

先掌握最小函数、简单循环、数组访问、函数调用、ABI 和测量误差，再进入后续卷的现代 CPU 性能、SIMD、HPC 和 AI 算子。

\section*{学习节奏建议}

第一遍学习时，建议按下面节奏推进：

\begin{itemize}
  \item 第 1-5 章不要急着优化代码，重点是跑通工具链、看懂地址、寄存器、对象布局和调试输出。
  \item 第 6-10 章每天都要读反汇编。只读汇编语法不够，必须把每个例子从 C++ 函数签名追到参数寄存器、栈和返回值。
  \item 第 11-12 章不要追求漂亮数字，重点是建立可复现实验结构：输入、热路径、校验、重复样本、CSV、环境记录和报告。
\end{itemize}

每章至少做一个实验、一个反例和一份短报告。实验验证模型，反例暴露误解，报告训练表达。对底层教材来说，这三者缺一项，理解都会不稳。

\section*{怎样使用 Linux 源码案例}

本书附录提供 Linux 源码阅读案例。它们不是要求你一开始就读完整内核，而是给每个系统现象一个真实源码入口。阅读这些案例时，先用用户态小程序制造现象，再用 \code{strace}、\code{readelf}、\filepath{/proc}、GDB 或 \code{perf} 取得工具证据，最后进入内核源码寻找对应路径。

源码阅读必须记录版本。你读到的路径属于某个 Linux tag、commit 或发行版源码包；如果没有运行同一份内核，就不能把源码阅读直接写成当前机器行为的证明。正确写法是：当前工具输出证明了用户态现象，某版本 Linux 源码解释了这种机制通常如何实现，二者共同支持一个有边界的结论。

建议在读完第 4 章后开始做附录中的 \code{write}、ELF 加载和 \filepath{/proc/self/maps} 案例；读完第 11 章后再做缺页、调度噪声和 \code{perf_event_open} 案例。不要急着扩大范围，先把一条窄路径读清楚。

\section*{遇到难点时的处理方式}

如果某一章读不动，通常不是因为某个术语太难，而是前面某一层没有闭合。可以按下面顺序排查：

\begin{enumerate}
  \item 看不懂汇编，先回到 C++ 函数签名和 ABI 参数位置。
  \item 看不懂地址表达式，先回到 \code{sizeof}、数组连续性和指针加法。
  \item 看不懂栈变化，先画 \register{rsp}、返回地址、参数和局部对象。
  \item 看不懂 benchmark 结果，先确认热路径有没有被编译器删除或内联改写。
  \item 看不懂性能差异，先确认输入规模、访问模式、cache 状态和测量重复次数。
\end{enumerate}

不要用背术语替代理解。如果一个概念说不清，就写一个 20 行以内的最小程序，用编译命令、反汇编和 GDB 把它拆开。

\section*{报告要求}

每个性能实验报告至少包含：

\begin{itemize}
  \item 实验目标。
  \item 机器和环境。
  \item 编译器和编译参数。
  \item correctness reference。
  \item 输入规模。
  \item benchmark 方法。
  \item 汇编、IR 或 profiling 证据。
  \item 结果表格。
  \item 失败实验。
  \item 结论和下一步验证。
\end{itemize}

\section*{报告的最低质量线}

一份合格报告必须能被别人复现。它至少要回答：

\begin{itemize}
  \item 代码从哪里来，提交或文件路径是什么。
  \item 用什么命令编译，优化级别和 ISA 参数是什么。
  \item 在什么机器和系统上运行。
  \item 输入如何生成，规模和分布是什么。
  \item 测的是哪一段代码，哪些工作不在计时区域内。
  \item 结果是否正确，正确性 reference 是什么。
  \item 原始样本在哪里，汇总口径是什么。
  \item 反汇编或工具输出是否支持你的解释。
  \item 结论有哪些限制，下一步需要验证什么。
\end{itemize}

写报告时宁可保守，不要把一个实验夸大成普遍规律。高质量教材训练的不是“说得像专家”，而是“证据能支撑到哪里，就说到哪里”。
